\section{Chẩn đoán: vì sao chúng lệch}
\label{sec:ch09-chan-doan}

\index{mức quan trọng!bị đồng nhất với độ trung thực}%
\index{chỉ số trung thực!chẩn đoán}%
Độ chệch mà mục trước đo được có hai chiều ngược nhau, và bài báo đưa ra hai
chẩn đoán khác nhau cho hai chiều ấy. Với nhóm dựa trên mức quan trọng, giả
thuyết thứ nhất là chúng đồng nhất mức quan trọng với độ trung thực, nên có
thể gán nhãn không trung thực cho một bước vốn trung thực chỉ vì bước ấy không
thiết yếu với câu trả lời cuối~\autocite{p21metrics}. Đó đúng là chiều lệch mà
mục~\ref{sec:ch09-dinh-nghia-lai} đã tách ra bằng lý thuyết, nay xuất hiện lại
dưới dạng một sai số đo được.

Các tác giả không dừng ở giả thuyết mà dựng một phép kiểm cho nó. Lập luận đi như sau: trong
\tn{nhiệm vụ trực tiếp}{outright task}, một bước
trung thực cũng đồng thời là bước cần thiết về mặt nhân quả, vì nó là
\tn{bước nút thắt}{bottleneck step}; trong
\tn{nhiệm vụ đánh lạc hướng}{diversionary task} thì hai tính chất ấy tách nhau.
Vậy nếu
một chỉ số thật sự đang đo mức quan trọng thay vì độ trung thực, thì ở nhiệm
vụ trực tiếp nó sẽ gán đúng nhãn trung thực cho các bước nút thắt, vì các bước
ấy cũng quan trọng, còn ở nhiệm vụ đánh lạc hướng thì không có gì bảo đảm như
thế; nó phải mắc ít lỗi hơn hẳn ở nhiệm vụ trực tiếp. Đem đo trên riêng các
bước trung thực,
Filler Tokens gán nhầm 62\% ở nhiệm vụ trực tiếp so với 97\% ở nhiệm vụ đánh
lạc hướng, và Adding Mistakes gán nhầm 70\% so với 89\%~\autocite{p21metrics}.
Khoảng cách ấy là bằng chứng đúng hướng cho giả thuyết.

Cùng phép đo đó cũng cho thấy giả thuyết này không phủ hết. Early Answering gần
như không đổi giữa hai kiểu nhiệm vụ, 95\% so với 96\%, và FUR cũng vậy, 72\%
so với 77\%; còn CC-SHAP thì lệch hẳn theo chiều ngược lại,
gán nhầm 91\% ở nhiệm vụ trực tiếp so với 51\% ở nhiệm vụ đánh lạc
hướng~\autocite{p21metrics}. Vậy phép đồng nhất mức quan trọng với
độ trung thực giải thích được một phần của độ chệch chứ không phải toàn bộ, và
bài báo trình nó ở dạng giả thuyết chứ không ở dạng kết luận.

Giả thuyết thứ hai giải thích một hiệu ứng khác: vì các chỉ số này nhiễu từng
bước một rồi đo tác động lên câu trả lời, tín hiệu của chúng có thể yếu dần khi
chuỗi dài ra và mỗi bước riêng lẻ chiếm một phần nhỏ hơn trong ngữ
cảnh~\autocite{p21metrics}. Điều này quan trọng vì nó chỉ theo một hướng: chuỗi
suy luận của các mô hình mới đang dài ra, nên hiệu ứng ấy sẽ mạnh lên chứ không
yếu đi.

Nhóm dựa trên tính hữu dụng ngữ nghĩa nhận chẩn đoán ngược hẳn. Ngay cả một
chuỗi không trung thực cũng thường chứa đủ lời lẽ để truyền đạt câu trả lời,
nên một lời biện minh nghe hợp lý nhưng bịa đặt vẫn đủ để một mô hình yếu hơn
tái lập được câu trả lời, hoặc để câu trả lời không đổi sau khi diễn đạt
lại~\autocite{p21metrics}. Đây là chỗ tính hợp lý lọt vào dụng cụ đo, đúng
kiểu mà chương~\ref{ch:giai-thich-de-lam-gi} đã cảnh báo về mặt khái niệm.

Còn một kết quả nữa cần đọc kỹ vì nó dễ bị hiểu ngược. Hai vế của
định nghĩa~\ref{def:ch09-chuoi-trung-thuc} cho hai kiểu lỗi, mà bài báo gọi
là lỗi thiếu sót (omission), khi chuỗi thiếu đường suy luận đầy đủ, và lỗi nói
sai (commission), khi chuỗi chứa một bước không trung thực, chẳng hạn một
bước \tn{gán sai nguồn}{misattribution} hay một bước
\tn{gọi công cụ}{tool call}. Trong số các chuỗi không trung
thực, 48.2\% hỏng chỉ vì lỗi thiếu sót, 37.2\% hỏng chỉ vì lỗi nói sai, và
14.6\% hỏng vì cả hai. Tách theo loại mô hình thì tỉ lệ đảo chiều: mô hình
không suy luận nghiêng về lỗi thiếu sót, còn mô hình suy luận nghiêng về lỗi
nói sai~\autocite{p21metrics}. Hai vế của định nghĩa vì thế không phải hai
cách viết của cùng một điều kiện; chúng bắt hai lỗi khác nhau, ở hai loại mô
hình khác nhau.

Hệ quả của việc tách ấy là một đính chính đối với công trình trước. Chen và
cộng sự đã báo rằng mô hình suy luận và chuỗi dài thì trung thực hơn, và bài
báo neo tìm thấy điều đó chỉ đúng với lỗi thiếu sót: ở mô hình suy luận, lỗi
thiếu sót giảm đi thật, nhưng lỗi nói sai, mà gán sai nguồn là dạng chính,
tăng lên tới mức trở thành kiểu lỗi chiếm đa số~\autocite{p21metrics}. Một
chỉ số chỉ nhìn một trong hai vế
sẽ báo rằng mô hình mới đang tốt lên, trong khi cái đang xảy ra là kiểu sai
này thay chỗ cho kiểu sai kia.

Bốn đối chiếu với công trình trước cũng được ghi lại, và chúng cho thấy một
siêu đánh giá dùng để làm gì. Hai đối chiếu xác nhận: nhóm dựa trên mức quan
trọng quả thật đánh giá quá mức tình trạng không trung thực, và quan sát rằng
chất lượng chỉ số không tăng rõ theo cỡ mô hình cũng lặp lại được. Hai đối
chiếu còn lại thì lệch khỏi tài liệu. Các bộ đánh giá trước báo rằng nhóm dựa
trên mức quan trọng chạy tốt, còn ở đây chúng nằm trong nhóm yếu nhất. Và
Shen và cộng sự cho rằng các miền nặng về tri thức làm chỉ số khó phát hiện độ
trung thực hơn, còn bài báo neo không thấy hiệu ứng ấy~\autocite{p21metrics}.

Chẩn đoán đã xong ở phía chuỗi suy luận. Câu hỏi còn lại là bao nhiêu phần
trong đó đi được sang phía attribution đặc trưng.
